\documentclass[11pt,a4paper]{amsart}

\usepackage[utf8]{inputenc}
\usepackage[english]{babel}
\usepackage{amsmath,amssymb,amsthm,mathtools}
\usepackage{microtype}
\usepackage{enumitem}
\usepackage{booktabs}
\usepackage{xcolor}
\usepackage[a4paper,margin=30mm]{geometry}
\usepackage[colorlinks=true,linkcolor=blue!55!black,citecolor=green!45!black,urlcolor=blue!60!black]{hyperref}

\newtheorem{theorem}{Theorem}[section]
\newtheorem{proposition}[theorem]{Proposition}
\newtheorem{lemma}[theorem]{Lemma}
\newtheorem{corollary}[theorem]{Corollary}
\theoremstyle{definition}
\newtheorem{definition}[theorem]{Definition}
\newtheorem{remark}[theorem]{Remark}

\newcommand{\e}{\varepsilon}
\newcommand{\dd}{\,\mathrm d}
\newcommand{\1}{\mathbf 1}
\newcommand{\cR}{\mathcal R}
\newcommand{\cU}{\mathcal U}
\newcommand{\cV}{\mathcal V}
\newcommand{\cM}{\mathcal M}
\newcommand{\Etwo}{E_2}
\newcommand{\lean}[1]{\texttt{#1}}

\title[Exact semiprimes in shorter logarithmic intervals]
{Exact semiprimes in almost all intervals\\
 of length \texorpdfstring{$(\log x)^{2.092}$}{(log x)\^{}2.092}}
\date{3 October 2026. Research draft; not peer reviewed.}

\begin{document}
\emergencystretch=2em

\begin{abstract}
\sloppy
Matom\"aki and Ter\"av\"ainen proved that almost every interval
\((x,x+(\log x)^{2.1}]\) contains a product of exactly two primes.  We
show that the exponent \(2.1\) can be replaced by any fixed
\[
 c>\frac{3688}{1763}=2.091888\ldots,
\]
and in particular by \(2.092\); quantitatively, almost every such
interval contains \(\gg(\log x)^{1.092}\) products \(pq\) with
\((\log x)^{1.091}<p\le(\log x)^{1.092}\).  This beats the
Matom\"aki--Ter\"av\"ainen exponent \(2.1\) by more than \(0.008\), and
it also goes below \(197/94\), the limit of their argument as it stands,
and below \(67/32\), the value obtained by inserting the Guth--Maynard
large-value theorem with a single threshold.  The proof keeps the
Harman-sieve and sparse-mean-value architecture of Matom\"aki and
Ter\"av\"ainen and changes how the large values of the two type~II
factors are controlled.  Instead of one large-value threshold for both
factors, we use the Guth--Maynard large-value estimates with a threshold
that depends on the length of the polynomial: the short factor is
powered into the range \([X^{5/6},X)\), where the Guth--Maynard
long-polynomial bound allows every exponent below \(1/4\), and the long
factor of length \(X^{1-\theta}\) gets the threshold
\(\frac{3-8\theta}{10(1-2\theta)}\) (for \(1/6\le\theta\le2/11\); it is
\(1/4\) for smaller \(\theta\)).  A three-case analysis of the
medium large-value bins then closes the two-branch sparse argument for
every \(a>1925/1763\), where \(P_1=(\log X)^a\).  The whole deduction,
from the cited published theorems and five explicitly labelled
uniform-in-parameter versions of Matom\"aki--Ter\"av\"ainen lemmas, has
been checked in Lean~4 with Mathlib.
\end{abstract}

\maketitle

\section{Introduction}

Let
\[
 \Etwo:=\{pq:p,q\text{ prime}\}
\]
be the set of integers with exactly two prime factors, counted with
multiplicity.  Unlike integers with \emph{at most} two prime factors,
\(\Etwo\) is subject to the parity problem, and finding its elements in
short intervals requires genuinely bilinear (type~II) information.

The previous record is due to Matom\"aki and Ter\"av\"ainen.

\begin{theorem}[{Matom\"aki--Ter\"av\"ainen \cite[Theorem~1.1]{MT}}]
\label{thm:MT}
There exist constants \(c_0>0\) and \(\delta>0\) such that, for all
\(X\ge3\), all but \(O(X/(\log X)^\delta)\) integers \(x\in[2,X]\)
satisfy
\[
 \#\bigl\{p_1p_2\in(x,x+(\log x)^{2.1}]:
 (\log x)^{1.09}<p_1\le(\log x)^{1.1}\bigr\}
 \ge c_0(\log x)^{1.1}.
\]
\end{theorem}

Our main result shortens the interval.

\begin{theorem}[Exponent \(2.092\)]\label{thm:clean}
There are constants \(c_0,\delta>0\) such that, for all sufficiently
large \(X\), all but \(O(X/(\log X)^\delta)\) integers \(x\le X\)
satisfy
\[
 \#\bigl\{pq\in(x,x+(\log x)^{2.092}]:
 p,q\text{ prime},\ (\log x)^{1.091}<p\le(\log x)^{1.092}\bigr\}
 \ge c_0(\log x)^{1.092}.
\]
In particular, almost every interval \((x,x+(\log x)^{2.092}]\)
contains an element of \(\Etwo\).
\end{theorem}

The exponent \(2.092\) is a convenient decimal; the method gives the
following.

\begin{theorem}[Parameter form]\label{thm:parameter}
For every fixed
\[
 c>\frac{3688}{1763}=1+\frac{1925}{1763}=2.0918888\ldots,
\]
almost every interval \((x,x+(\log x)^c]\) contains an element of
\(\Etwo\); the number of exceptional integers \(x\le X\) is
\(O_c(X/(\log X)^{\delta_c})\) for some \(\delta_c>0\).  If moreover
\(c\le21/10\) and \(\Delta>0\) satisfies \(c-1-\Delta>1925/1763\), then,
outside an exceptional set of the same type, the interval
\((x,x+(\log x)^c]\) contains \(\gg_{c,\Delta}(\log x)^{c-1}\) products
\(pq\) with \((\log x)^{c-1-\Delta}<p\le(\log x)^{c-1-\Delta/2}\).
\end{theorem}

\subsection*{How this compares with Matom\"aki--Ter\"av\"ainen}
Theorem~\ref{thm:clean} beats the published exponent \(2.1\) of
\cite[Theorem~1.1]{MT}: the interval length drops by the factor
\[
 \frac{(\log x)^{2.1}}{(\log x)^{2.092}}=(\log x)^{0.008}\longrightarrow\infty ,
\]
and the threshold \(3688/1763\) is below \(2.1\) by
\[
 \frac{21}{10}-\frac{3688}{1763}=\frac{143}{17630}>0.0081 .
\]
The improvement is not an artefact of the decimals chosen in \cite{MT}.
Their argument, optimized as it stands (Jutila's large-value threshold
\(49/206\) in \cite[Lemma~4.1]{MT} and the two sparse branches of
\cite[Section~5.3]{MT}), reaches \(c>1+103/94=197/94\) and no further;
see the last display of \cite[Section~5.3]{MT}.  Replacing Jutila's
estimate by the Guth--Maynard large-value theorem with a single threshold
\(17/70\) for both type~II factors gives \(c>67/32\).  The present paper
goes below both:

\begin{center}
\small
\begin{tabular}{@{}lll@{}}
\toprule
Argument & large-value input & admissible \(c\)\\
\midrule
\cite[Thm.~1.1]{MT}, as printed & Jutila, \(49/206\) & \(c=2.1\)\\
\cite{MT}, optimized & Jutila, \(49/206\) & \(c>197/94=2.09574\ldots\)\\
single threshold & Guth--Maynard, \(17/70\) & \(c>67/32=2.09375\)\\
\textbf{this paper} & \textbf{Guth--Maynard, by length} & \(\boldsymbol{c>3688/1763=2.09188\ldots}\)\\
\bottomrule
\end{tabular}
\end{center}

\noindent In the language of \cite{MT}: a pair of polynomial lengths
\(M_1=X^\theta\), \(M_2=X^{1-\theta}\) with \(\theta\le2/11\) never
needs the worst threshold \(17/70\) on \emph{both} factors.  The short
factor can always be powered into \([X^{5/6},X)\), and the long factor
needs \(17/70\) only when \(\theta=2/11\) exactly.  Section~\ref{sec:bins}
turns this observation into the bound \(a>1925/1763\).

\subsection*{Formal verification}
The deduction of Theorems~\ref{thm:clean} and~\ref{thm:parameter} has
been formalized in Lean~4 with Mathlib (Section~\ref{sec:lean}).  The
Lean theorems take two explicit hypothesis bundles: published theorems
transcribed in their printed ranges (Guth--Maynard, Heath-Brown,
Hildebrand--Tenenbaum, Matom\"aki--Ter\"av\"ainen, Iwaniec--Kowalski,
\dots), and five labelled uniform-in-parameter versions of
Matom\"aki--Radziwi\l\l, Ter\"av\"ainen and Matom\"aki--Ter\"av\"ainen
lemmas that are printed only for \(c=2.1\).  All of the new mathematics
of this paper (Lemma~\ref{lem:refined}, Proposition~\ref{prop:bins} and
everything after) is proved in Lean, not assumed.

\subsection*{Status and provenance}
This is a research draft and has not been refereed.  An earlier draft,
generated with ChatGPT (OpenAI), obtained \(c>67/32\) with a single
threshold \(17/70\); the length-dependent refinement, the present
write-up and the Lean formalization of both were produced with Claude
(Anthropic).  Readers should check the argument independently before
relying on it.

\section{Notation}\label{sec:notation}

All logarithms are natural; \(p,q\) denote primes and \(n\sim N\) means
\(N<n\le2N\).  We write \(A\ll B\) or \(A=O(B)\) if \(|A|\le CB\) for a
constant \(C\) depending only on the displayed subscripts and on fixed
parameters.  The parameters are chosen in the following order: first the
gap \(\nu>0\) between \(a\) and \(1925/1763\); then the analytic
parameter \(\e>0\), small in terms of \(\nu\) and of the absolute
constant \(C_1\) of Proposition~\ref{prop:typeII}; finally \(X\) is
taken large in terms of everything else.

A sequence \((a_n)\) is \((A_0,B)\)-\emph{divisor bounded} if
\(|a_n|\le A_0d(n)^B\) on its support, with \(A_0\ge1\), \(B\ge0\)
independent of the principal scales.  If \(a\) is \((A_0,B)\)-divisor
bounded, then the coefficients of the \(\ell\)-fold Dirichlet
convolution of \(a\) restricted to a dyadic range are
\((A_0^\ell,B\ell+\ell-1)\)-divisor bounded: the first \(\ell-1\)
factors of an ordered factorization of \(n\) are divisors of \(n\).

For a Dirichlet polynomial \(A\), a height \(T\) and a level \(V>0\) we
write
\[
 \mathcal L(A;T,V)=\{t\in[0,T]:|A(1+it)|>V\}.
\]

\section{The framework of Matom\"aki and Ter\"av\"ainen}
\label{sec:framework}

We recall the part of \cite{MT} that we use, with the parameters left
free.  Let \(c\in(2,21/10]\), \(h=(\log X)^c\), and let
\(P_1=(\log X)^a\) with \(1<a\le c-1\le 11/10\).  Harman's sieve with
\(z=X^{2/11}\) provides a minorant \(\rho^-(n)\le\1_{\mathbb P}(n)\)
\cite[(2.3)]{MT} with positive average on long intervals
\cite[Theorem~2.1(ii)]{MT}.  By Perron's formula
\cite[Lemma~3.1]{MT} (after \cite[Lemma~14]{MR} and \cite[Lemma~1]{Ter}),
the variance of \(\sum_{x<pn\le x+h,\,p\sim P_1}\rho^-(n)\) is controlled
by mean values
\[
 \int_{X^{1/1000}}^{T}|P_1(1+it)|^2|P(1+it)|^2\dd t,\qquad
 P_1(s)=\sum_{p\sim P_1}p^{-s},\quad
 P(s)=\sum_{X/(2P_1)<n\le4X/P_1}\frac{\rho^-(n)}{n^s}.
\]
On the set
\begin{equation}\label{eq:Udef}
 \cU=\bigl\{t\in[X^{1/1000},X]:|P_1(1+it)|\ge P_1^{-\e/10}\bigr\}
\end{equation}
the polynomial \(P\) is decomposed by \cite[Proposition~2.2]{MT} into
\(O(\exp((\log\log X)^5))\) type~I, type~I/II and type~II convolutions
and a negligible remainder.  The type~I and type~I/II components are
estimated in \cite[Sections~5.1--5.2]{MT}; their only dependence on
\(a\) is through the power savings
\(X^{-(1-1/a)/20+O(\e)}\) and \(X^{-\e(1-1/a)/4+\e/100+O(\e^2)}\), which
remain power savings for \(a\ge1925/1763\) because
\[
 \frac{1-1/a}{4}\ge\frac{1-1763/1925}{4}=\frac{81}{3850}>\frac1{100}.
\]

A type~II component is \(F=M_1M_2\), where
\[
 M_1(s)=\sum_{m\sim M_1}\frac{\alpha_m}{m^s},\qquad
 M_2(s)=\sum_{n\sim M_2}\frac{\beta_n}{n^s},\qquad
 M_1M_2\asymp\frac{X}{P_1},\qquad X^{\e/2}\le M_1\le X^{2/11},
\]
with divisor-bounded coefficients and \(\alpha\) of the structured form
\cite[(2.11)]{MT}.  We write \(\theta=\log M_1/\log X\in[\e/2,2/11]\);
then \(M_2=X^{1-\theta-o(1)}\).  The target is
\begin{equation}\label{eq:component-target}
 \int_{\cU}|P_1(1+it)|^2|M_1(1+it)M_2(1+it)|^2\dd t
 \ll\exp\bigl(-(\log\log X)^6\bigr).
\end{equation}
As in \cite[Section~5.3]{MT}, \(\cU\) is split into the set \(\cU_1\)
where \(|M_j(1+it)|\ge M_j^{-\sigma_*}\) for some \(j\), the set
\(\cU_3\) where some \(|M_j(1+it)|<X^{-1}\), and the rest \(\cU_2\),
which is cut into \(O((\log X)^2)\) bins
\begin{equation}\label{eq:Vdef}
 \cV_{\sigma_1,\sigma_2}=\bigl\{t\in\cU_2:
 M_j^{-\sigma_j}<|M_j(1+it)|\le eM_j^{-\sigma_j}\ (j=1,2)\bigr\},
 \qquad \sigma_1,\sigma_2>\sigma_* .
\end{equation}
The set \(\cU_1\) is handled by the type~II estimate
(Proposition~\ref{prop:typeII}) together with the Vinogradov--Korobov
bound \cite[(5.6)]{MT}
\(\sup_{t\in\cU}|M_1(1+it)|\ll\exp(-2(\log X)^{1/10})\); \(\cU_3\) is
trivial.  Everything therefore reduces to the bin estimate
\begin{equation}\label{eq:Vtarget}
 |\cV_{\sigma_1,\sigma_2}|\ll M_1^{2\sigma_1}M_2^{2\sigma_2}X^{-\e^2/4},
\end{equation}
uniformly in \(\sigma_1,\sigma_2>\sigma_*\) and in the component.  This
is where the exponent \(a\), and hence \(c\), is decided.

\section{Length-dependent large-value thresholds}\label{sec:largevalues}

We use two results of Guth and Maynard.

\begin{theorem}[{Guth--Maynard \cite[Theorem~1.1]{GM}}]\label{thm:GM}
Let \(|b_n|\le1\) and let \(t_1,\dots,t_R\in[0,T]\) be one-spaced with
\(|\sum_{N<n\le2N}b_nn^{it_r}|\ge V\) for all \(r\).  Then
\[
 R\le T^{o(1)}\bigl(N^2V^{-2}+N^{18/5}V^{-4}+TN^{12/5}V^{-4}\bigr).
\]
\end{theorem}

\begin{theorem}[{Guth--Maynard \cite[Proposition~12.1]{GM}}]
\label{thm:GM-long}
Under the hypotheses of Theorem~\ref{thm:GM}, if
\(T^{5/6}\le N\le T\) and \(V=N^{\varrho}\) with \(\varrho\ge7/10\),
then
\[
 R\le T^{o(1)}\bigl(N^{2-2\varrho}+T^{1/2}N^{3-4\varrho}
 +T^{(30\varrho-21)/5}N^{(46-60\varrho)/5}\bigr).
\]
\end{theorem}

For \(9/11\le\eta\le1\) put
\[
 s(\eta)=
 \begin{cases}
 \dfrac{8\eta-5}{10(2\eta-1)}, & 9/11\le\eta\le5/6,\\[2mm]
 \dfrac14, & 5/6\le\eta\le1,
 \end{cases}
 \qquad
 s\Bigl(\frac9{11}\Bigr)=\frac{17}{70},\quad s\Bigl(\frac56\Bigr)=\frac14 .
\]
The function \(s\) is increasing and continuous.  The value \(17/70\) used
in the single-threshold argument is its minimum, attained only at the
shortest length \(T^{9/11}\).

\begin{lemma}[Refined density estimate]\label{lem:refined}
Let \(0<\e\le1/1000\), \(A_0\ge1\), \(B\ge0\).  There are \(C,T_0\)
such that the following holds for \(T\ge T_0\).  Let
\(\eta_0\in[9/11,5/6]\), let \(N\) be an integer with
\(T^{\eta_0-10\e}\le N\le T^{1-\e/10}\), let \(A(s)=\sum_{n\sim N}a_nn^{-s}\)
with \((A_0,B)\)-divisor bounded coefficients, and let \(\sigma\) satisfy
\begin{equation}\label{eq:sigma-range}
 10\e\le\sigma\le\frac14-20\e,\qquad
 \sigma(4\eta_0-2)+20\e\le\frac{8\eta_0}{5}-1 .
\end{equation}
Then
\[
 \bigl|\mathcal L(A;T,N^{-\sigma})\bigr|\le C\,T^{2\sigma-\e^2/2}.
\]
\end{lemma}

The second condition in \eqref{eq:sigma-range} says
\(\sigma\le s(\eta_0)-O(\e)\) when \(\eta_0<5/6\); at \(\eta_0=5/6\) it
is implied by the first.

\begin{proof}
Write \(N=T^\eta\).  Normalize as in \cite[Lemma~4.1]{MT}: with a fixed
\(0<\varpi\le\e^2/1000\), choose \(D=D(B,\varpi)\) with \(d(n)^B\le Dn^\varpi\)
and put \(b_n=\frac{N}{n}\,a_n/(A_0D2^\varpi N^\varpi)\), so that
\(|b_n|\le1\) and \(|A(1+it)|>N^{-\sigma}\) implies
\(|\sum b_nn^{-it}|\ge N^{1-\sigma-2\varpi}\) for large \(T\).  Pass from
measure to one-spaced sets by the standard local discretization, and
reflect \(t\mapsto T-t\) if needed.

\emph{Case \(\eta\le5/6\).}  Theorem~\ref{thm:GM} with
\(V=N^{1-\sigma-2\varpi}\) gives
\[
 R\le T^{o(1)}\bigl(N^{2\sigma+4\varpi}+N^{4\sigma-2/5+8\varpi}
 +TN^{4\sigma-8/5+8\varpi}\bigr).
\]
Compare each exponent of \(T\) with \(2\sigma\).  The first is
\(2\sigma\eta\le2\sigma-2\sigma\e/10\le2\sigma-2\e^2\).  For the second,
since \(4\eta-2>0\) and \(\sigma\le1/4-20\e\),
\[
 \eta(4\sigma-\tfrac25)-2\sigma=\sigma(4\eta-2)-\tfrac{2\eta}5
 \le\tfrac{3\eta}5-\tfrac12-20\e(4\eta-2)<-20\e,
\]
using \(\eta\le5/6\).  The third exponent
\(1+\eta(4\sigma-8/5)\) is decreasing in \(\eta\), so it is at most its
value at \(\eta=\eta_0-10\e\):
\[
 1+\eta_0\bigl(4\sigma-\tfrac85\bigr)+10\e\bigl(\tfrac85-4\sigma\bigr)
 \le 2\sigma-20\e+16\e=2\sigma-4\e
\]
by \eqref{eq:sigma-range}.

\emph{Case \(\eta\ge5/6\).}  Theorem~\ref{thm:GM-long} with
\(\varrho=1-\sigma-2\varpi\ge7/10\) gives the exponents
\(2\sigma\eta+O(\varpi)\),
\[
 \tfrac12+\eta(4\sigma-1)+O(\varpi)
 =2\sigma-(2\eta-1)\bigl(\tfrac12-2\sigma\bigr)+O(\varpi)
 \le2\sigma-\tfrac{80}3\e+O(\varpi),
\]
and
\begin{align*}
 \tfrac{9-30\sigma}5+\eta\,\tfrac{60\sigma-14}5+O(\varpi)
 &=2\sigma+\tfrac15\bigl(9-14\eta-\sigma(40-60\eta)\bigr)+O(\varpi)\\
 &\le2\sigma-\tfrac15\bigl(1-\eta+20\e(60\eta-40)\bigr)+O(\varpi),
\end{align*}
where we used \(\sigma\le1/4-20\e\) and \(60\eta-40>0\).  All savings
are at least \(2\e^2-O(\varpi)\), and the \(T^{o(1)}\) and normalization
losses are absorbed for large \(T\).
\end{proof}

\begin{remark}
Lemma~\ref{lem:refined} is exactly what the two Guth--Maynard bounds
give; nothing beyond \cite{GM} is used.  The single-threshold argument
uses only the value \(s(9/11)=17/70\) and applies it to both type~II
factors, whatever their lengths.
\end{remark}

\section{The type~II estimate with threshold \texorpdfstring{\(17/70\)}{17/70}}
\label{sec:typeII}

Lemma~\ref{lem:refined} with \(\eta_0=9/11\) and \(\sigma\le17/70-O(\e)\)
covers every length in \([T^{9/11-10\e},T^{1-\e/10}]\).  Inserting it
into the proof of \cite[Proposition~4.2]{MT} in place of
\cite[Lemma~4.1]{MT} gives the following.

\begin{proposition}[Type~II, threshold \(17/70\)]\label{prop:typeII}
There is an absolute constant \(C_1\) such that the following holds for
\(\e>0\) small.  Let \(M_1,M_2\) be Dirichlet polynomials of lengths
\(M_1,M_2\) with divisor-bounded coefficients,
\((\log T)^{-L}\le M_1M_2/T\le(\log T)^L\) and
\(T^{\e/5}\le M_1\le T^{2/11+\e}\), and put
\(\sigma_*=17/70-C_1\e\).  If \(\cU'\subset[0,T]\) is measurable and on
\(\cU'\) one has \(|M_1(1+it)|\ge M_1^{-\sigma_*}\) or
\(|M_2(1+it)|\ge M_2^{-\sigma_*}\), then
\[
 \int_{\cU'}|M_1(1+it)M_2(1+it)|^2\dd t
 \ll T^{-\kappa_1}+(\log T)^{C_2}\sup_{t\in\cU'}|M_1(1+it)|^2
\]
for some \(\kappa_1=\kappa_1(\e,B)>0\) and \(C_2=C_2(\e,A_0,B,L)\).
\end{proposition}

\begin{proof}
As in \cite[Proposition~4.2]{MT}: the large-value parts \(\cR_1\),
\(\cR_3\) are unchanged; in the middle range one powers \(M_1\) into
\([T^{5/6-\e},T^{1-\e/2}]\) when \(\tau_1\le\tau_2\), or applies the
density bound to \(M_2\) (of length in \([T^{9/11-2\e},T^{1-\e/6}]\))
when \(\tau_2<\tau_1\), now with Lemma~\ref{lem:refined} at
\(\eta_0=9/11\) instead of Jutila's estimate.
\end{proof}

This is the large-value set \(\cU_1\) of Section~\ref{sec:framework}.
The new work is in the bins \(\cV_{\sigma_1,\sigma_2}\), where both
\(\sigma_j>\sigma_*\).

\section{The medium bins}\label{sec:bins}

Fix \(\nu>0\) and let
\[
 \frac{1925}{1763}+\nu\le a\le\frac{11}{10},\qquad u=1-\frac1a .
\]
Then \(u\ge162/1925+\frac{10}{11}\cdot\frac{1763}{1925}\nu\).  Consider
one bin \(\cV=\cV_{\sigma_1,\sigma_2}\) with \(\sigma_1,\sigma_2>\sigma_*\),
and write
\[
 \theta=\frac{\log M_1}{\log X}\in\Bigl[\frac\e2,\frac2{11}\Bigr],
 \qquad
 \phi=\frac{\log M_2}{\log X},\qquad
 1-\theta-o(1)\le\phi\le1-\theta .
\]

\subsection{The sparse branches}
We recall \cite[Section~5.3]{MT}.  With
\(k=\lfloor\log X/\log(2P_1)\rfloor\) and \(M=P_1^k=X^{1-o(1)}\), the
sparse polynomial \(Q=P_1^k/k!\) has
\(|Q(1+it)|\ge M^{-\e/10-1/a+o(1)}\) on \(\cU\) and support of size
\(M^{1-1/a+o(1)}\) by Hildebrand--Tenenbaum \cite{HT}.  Heath-Brown's
sparse mean value theorem \cite[Theorem~4]{HB}, applied to \(Q\) and
\(M_2\), proves \eqref{eq:Vtarget} unless
\begin{equation}\label{eq:residual}
 \theta(1-2\sigma_1)\ge u-O(\e)
 \qquad\text{(the first branch fails).}
\end{equation}
Applied to \(Q\) and \(M_1^5\) (the first term of \cite[Lemma~3.4]{MT} is
absorbed because \(M_1^5\le X^{10/11}\le X^{1/a}\)), it proves
\eqref{eq:Vtarget} whenever
\begin{equation}\label{eq:key}
 K(\theta,\sigma_1,\sigma_2):=(5-8\sigma_1)\theta+2\sigma_2(1-\theta)
 >\frac1a+O(\e).
\end{equation}
Note the identity
\begin{equation}\label{eq:identity}
 K(\theta,\sigma_1,\sigma_2)=4\theta(1-2\sigma_1)+\theta(1-2\sigma_2)+2\sigma_2 .
\end{equation}
Under \eqref{eq:residual} we have \(\theta\ge1/7\), since
\(1-2\sigma_1<1-2\sigma_*\le18/35+O(\e)\) and \(u\ge162/1925\).

\subsection{The density cases}
From now on assume that the first branch fails, so that \eqref{eq:residual}
holds and \(1/7\le\theta\le2/11\).
\begin{enumerate}[label=(\alph*),leftmargin=*]
\item \emph{\(\sigma_1\le\sigma_2\) and \(\sigma_1<1/4-O(\e)\).}  Take
\(\ell=6\) if \(6\theta\le1-\e/5\) and \(\ell=5\) otherwise; since
\(1/7\le\theta\le2/11\), the \(\ell\) dyadic pieces of \(M_1^\ell\) have
lengths in \([M_1^\ell,2^\ell M_1^\ell]\subset[X^{5/6-O(\e)},X^{1-\e/10}]\).
On \(\cV\), \(|M_1^\ell|\) exceeds \((M_1^\ell)^{-\sigma_1}\), so one of the
pieces exceeds \((M_1^\ell)^{-\sigma_1-o(1)}\).  Lemma~\ref{lem:refined}
with \(\eta_0=5/6\) and \(3\e/2\) in place of \(\e\) gives
\(|\cV|\ll X^{2\sigma_1-9\e^2/16}\le M_1^{2\sigma_1}M_2^{2\sigma_2}X^{-\e^2/4}\),
because \(M_1^{2\sigma_1}M_2^{2\sigma_2}\ge(M_1M_2)^{2\sigma_1}
=X^{2\sigma_1-o(1)}\).
\item \emph{\(\sigma_2<\sigma_1\) and \(\sigma_2<s(1-\theta)-O(\e)\).}
Apply Lemma~\ref{lem:refined} (with \(3\e/2\) in place of \(\e\)) directly
to \(M_2\), whose length is \(X^\phi\), with
\(\eta_0=\max(9/11,\min(5/6,1-\theta))\); in terms of \(\theta\) the
threshold is
\[
 s_2(\theta):=s(1-\theta)=
 \begin{cases}
 1/4, & \theta\le1/6,\\[1mm]
 \dfrac{3-8\theta}{10(1-2\theta)}, & 1/6\le\theta\le2/11,
 \end{cases}
 \qquad s_2\Bigl(\frac2{11}\Bigr)=\frac{17}{70}.
\]
This gives \(|\cV|\ll X^{2\sigma_2-9\e^2/8}\), which suffices as in (a).
\end{enumerate}

\subsection{The sparse cases}
If neither (a) nor (b) applies and the first branch fails, the second
branch closes.

\begin{proposition}\label{prop:bins}
Assume \eqref{eq:residual} and that neither (a) nor (b) applies.  Then
\eqref{eq:key} holds, provided \(\e\) is small in terms of \(\nu\).
\end{proposition}

\begin{proof}
There are three regions.

\emph{Region A: \(\sigma_1\le\sigma_2\), so \(\sigma_1\ge1/4-O(\e)\).}
By \eqref{eq:identity},
\(K=2\sigma_1+5\theta(1-2\sigma_1)+2(\sigma_2-\sigma_1)(1-\theta)
 \ge\frac12+5u-O(\e)\).

\emph{Region B1: \(\sigma_2<\sigma_1\) and \(\sigma_2\ge1/4-O(\e)\).}
Since \(\theta(1-2\sigma_2)\ge\theta(1-2\sigma_1)\), \eqref{eq:identity}
gives \(K\ge5\theta(1-2\sigma_1)+2\sigma_2\ge\frac12+5u-O(\e)\).

In both regions \(K>1/a=1-u\) as soon as \(6u>\frac12\), i.e.\
\(a>12/11\); and \(1925/1763>12/11\).

\emph{Region B2: \(\sigma_2<\sigma_1\), \(\sigma_2<1/4-O(\e)\) and
\(\sigma_2\ge s_2(\theta)-O(\e)\).}  Then \(\theta\ge1/6\) (otherwise
\(s_2(\theta)=1/4\)), and
\(\sigma_2(2-4\theta)\ge\frac35-\frac85\theta-O(\e)\).  By
\eqref{eq:identity} and \eqref{eq:residual},
\[
 K\ge4u+\theta+2\sigma_2(1-\theta)-O(\e).
\]
Multiplying by \(1-2\theta>0\),
\begin{align*}
 (1-2\theta)\bigl(\theta+2\sigma_2(1-\theta)\bigr)
 &\ge\theta(1-2\theta)+(1-\theta)\Bigl(\frac35-\frac85\theta\Bigr)-O(\e)\\
 &=\frac{223}{385}(1-2\theta)+\frac2{385}(2-11\theta)(7\theta+2)-O(\e),
\end{align*}
and \((2-11\theta)(7\theta+2)\ge0\) because \(\theta\le2/11\).  Hence
\(\theta+2\sigma_2(1-\theta)\ge\frac{223}{385}-O(\e)\) and
\[
 K\ge4u+\frac{223}{385}-O(\e)>1-u
 \iff 5u>\frac{162}{385}+O(\e)
 \iff \frac1a<\frac{1763}{1925}-O(\e),
\]
which holds since \(a\ge1925/1763+\nu\).
\end{proof}

Combining (a), (b), Proposition~\ref{prop:bins} and the first sparse
branch proves \eqref{eq:Vtarget} for every bin, uniformly for
\(a\in[1925/1763+\nu,11/10]\).

\begin{remark}[The extremal configuration]
The constant \(1925/1763\) is attained.  At
\[
 \theta=\frac2{11},\qquad\sigma_2=\frac{17}{70}=s_2\Bigl(\frac2{11}\Bigr),
 \qquad\sigma_1=\frac{517}{1925},\qquad a=\frac{1925}{1763},
\]
the first branch is exactly saturated, \(\theta(1-2\sigma_1)=\frac{162}{1925}=u\),
and so is the second:
\[
 (5-8\sigma_1)\theta+2\sigma_2(1-\theta)
 =\frac{998}{1925}+\frac{153}{385}=\frac{1763}{1925}=\frac1a .
\]
Neither density case applies there (\(\sigma_2<\sigma_1\) and \(\sigma_2\)
is exactly at the threshold of \(M_2\)).  If the long factor could also
use the threshold \(1/4\) at length \(X^{9/11}\), Region~B2 would
disappear and the argument would give \(a>12/11\), i.e.\ \(c>23/11\).
\end{remark}

\section{Completion of the proof}\label{sec:completion}

With \eqref{eq:Vtarget} in hand, the remainder is the argument of
\cite[Sections~2 and~5]{MT} with the parameters left free.

\subsection{The parameterized reduction}

\begin{proposition}\label{prop:variance}
Fix \(\nu>0\) and \(c\in(2,21/10]\), and let
\(1925/1763+\nu\le a\le\min(11/10,c-1)\), \(P_1=(\log X)^a\),
\(h=(\log X)^c\) and \(h_1=X^{99/100}\).  There is \(\e_0>0\),
depending only on \(\nu\) and \(c\), such that, uniformly in \(a\),
\[
 \frac1X\int_X^{2X}\Bigl|\frac1h\sum_{\substack{x<pn\le x+h\\p\sim P_1}}\rho^-(n)
 -\frac1{h_1}\sum_{\substack{x<pn\le x+h_1\\p\sim P_1}}\rho^-(n)\Bigr|^2\dd x
 \ll(\log X)^{-2-\e_0}.
\]
\end{proposition}

\begin{proof}
By the Perron/Parseval reduction \cite[Lemma~3.1]{MT} (whose proof is
uniform in \(a\) and \(c\)) it suffices to prove
\[
 \int_{X^{1/1000}}^T|P_1(1+it)P(1+it)|^2\dd t
 \ll\frac{T}{X/h}\,(\log X)^{-2-\e_0}\qquad(T\ge X/h).
\]
For \(T\ge X\) this follows from the mean value theorem and the fact that
each \(n\) has at most one prime factor \(p\sim P_1\).  On the complement
of \(\cU\) in \([X^{1/1000},X]\) one uses \(|P_1(1+it)|<P_1^{-\e/10}\),
the improved mean value theorem \cite[Lemma~3.3]{MT} and the
rough-number sieve bound of \cite[Section~5]{MT}.  On \(\cU\) one uses
the decomposition \cite[Proposition~2.2]{MT}.  The type~I and type~I/II
components are bounded as in \cite[Sections~5.1, 5.2]{MT}: the power
savings \(X^{-(1-1/a)/20+O(\e)}\) and
\(X^{-\e(1-1/a)/4+\e/100+O(\e^2)}\) left there are genuine because
\[
 \frac{1-1/a}{4}\ge\frac{1-1763/1925}{4}=\frac{81}{3850}>\frac1{100}.
\]
The type~II components are bounded by Proposition~\ref{prop:typeII} on
\(\cU_1\), trivially on \(\cU_3\), and by \eqref{eq:Vtarget},
Section~\ref{sec:bins}, on the \(O((\log X)^2)\) bins of \(\cU_2\).
\end{proof}

\subsection{From variance to exact semiprimes}

Fix \(c\in(3688/1763,21/10]\) and \(\Delta>0\) with
\(c-1-\Delta>1925/1763\), and take the disjoint dyadic prime blocks
\(p\sim P_1\) with
\((\log X)^{c-1-3\Delta/4}<P_1<2P_1\le(\log X)^{c-1-5\Delta/8}\).
For each block, \(a=\log P_1/\log\log X\) lies in
\([1925/1763+\nu,11/10]\) for some \(\nu=\nu(c,\Delta)>0\), so
Proposition~\ref{prop:variance} applies uniformly.  Summing over the
\(\asymp\log\log X\) blocks with Cauchy--Schwarz, the long-interval lower
bound \cite[Theorem~2.1(ii)]{MT}, Chebyshev's inequality and
\(\rho^-\le\1_{\mathbb P}\) give, for all \(x\in[X,2X]\) outside a set of
measure \(O(X/(\log X)^\delta)\),
\[
 \#\{pn\in(x,x+h]:p,n\text{ prime},\ p\text{ in a block}\}
 \gg_{c,\Delta}(\log X)^{c-1}.
\]
Three routine passages complete the proof of Theorem~\ref{thm:parameter}
for \(c\le21/10\): from real to integer \(x\) (using the same bound for
\(h/2\)); from \([X,2X]\) to \([1,X]\) (dyadic summation preserves the
exponent \(\delta\)); and from the blocks to the window
\[
 \bigl((\log x)^{c-1-\Delta},(\log x)^{c-1-\Delta/2}\bigr].
\]
For \(c>21/10\) the existence statement follows from the case
\(c=21/10\).

\subsection{The value \texorpdfstring{\(2.092\)}{2.092}}
Take \(c=2.092=523/250\) and \(\Delta=10^{-4}\).  Then
\[
 c-1-\Delta=1.0919=\frac{10919}{10000},\qquad
 \frac{10919}{10000}-\frac{1925}{1763}=\frac{197}{17630000}>0,
\]
so Theorem~\ref{thm:parameter} applies.  Its prime window satisfies
\[
 \bigl((\log x)^{1.0919},(\log x)^{1.09195}\bigr]
 \subset\bigl((\log x)^{1.091},(\log x)^{1.092}\bigr],
\]
and its lower bound is \(\gg(\log x)^{1.092}\).  This proves
Theorem~\ref{thm:clean}.

\section{The Lean formalization}\label{sec:lean}

The accompanying Lean~4 project (Lean~4.34.0, Mathlib v4.34.0) proves
\begin{verbatim}
theorem ExactSemiprimes.mainTheorem_exponent_2092
    (inputs : ExternalInputs) (ext : ParameterExtensionInputs) :
    Results.CleanTheoremStatement      -- Theorem 1.2
theorem ExactSemiprimes.mainTheorem_parameter
    (inputs : ExternalInputs) (ext : ParameterExtensionInputs) :
    Results.ParameterTheoremStatement  -- Theorem 1.3
\end{verbatim}
\begin{itemize}[leftmargin=*]
\item \lean{ExternalInputs} collects published theorems, each transcribed
  in its printed parameter range: Guth--Maynard's Theorem~1.1 and
  Proposition~12.1, Heath-Brown's sparse mean value theorem,
  Hildebrand--Tenenbaum's smooth-number bound, the decomposition,
  equations and lemmas of \cite{MT} that are used verbatim, the mean
  value theorem of Iwaniec--Kowalski, and a few others.
\item \lean{ParameterExtensionInputs} collects five labelled statements
  that are printed only at \(c=2.1\) (or \(a\in[1.0999,1.1]\)) and whose
  printed proofs are uniform in the parameters: the Parseval/Perron
  reduction \cite[Lemma~14]{MR}, \cite[Lemma~1]{Ter}, \cite[Lemma~3.1]{MT};
  the long-interval lower bound \cite[Theorem~2.1(ii)]{MT}; the rough-pair
  sieve bound of \cite[Section~5]{MT}; and the type~I and type~I/II
  bounds of \cite[Sections~5.1, 5.2]{MT} for
  \(a\in[1925/1763,11/10]\).
\item Everything else is proved: Lemma~\ref{lem:refined}
  (\lean{Hybrid/RefinedDensity.lean}), Proposition~\ref{prop:typeII},
  the two sparse branches, the three cases of Section~\ref{sec:bins}
  (\lean{Final/TypeIIDensityBins.lean}, \lean{Final/TypeIIComponent.lean}),
  the Perron target, the variance and counting arguments and the
  passage to Theorems~\ref{thm:clean} and~\ref{thm:parameter}.
\item There is no \lean{sorry} and no project \lean{axiom}.  The
  command \lean{\#print axioms} reports only the three standard axioms
  \lean{propext}, \lean{Classical.choice} and \lean{Quot.sound}.
\end{itemize}
The file \lean{ASSUMPTIONS.md} of the project lists every hypothesis
with its source.

\section*{Acknowledgements}

Siddharth Iyer was used as a ``meat proxy''\footnote{A \emph{meat proxy}
is a human who serves as the physical go-between for AI systems, much as a
network proxy forwards traffic: typing or pasting prompts in, carrying the
outputs from one model or tool to another, moving files around, and
running whatever the systems cannot run for themselves.} for this project.
He collected the outputs of the language models and told ChatGPT to
discover problems in the work of Matom\"aki.  The project began with the
following prompt, an important human idea, which he gave to ChatGPT:
\begin{quote}\small
Identify works of Matomaki and detail five problems from Matomaki's works
that seem within reach, they could be 1) an improvement of a result of
Matomaki, 2) Answering a question of Matomaki, 3) Disproving a question of
Matomaki.  For now you do not need to supply proofs for five such problems,
operate on vibes and slight rigour, also try and balance it so that the
problems are ``interesting'' to the broader mathematical community.  For
each such problem also detail a reference that relates to Matomakis work.
Detail the five problems in a md file, so that I can supply it towards
future agent runs (that actually attempt to solve some of the problems).
\end{quote}

\begin{thebibliography}{99}

\bibitem{GM}
L. Guth and J. Maynard,
\emph{New large value estimates for Dirichlet polynomials},
Ann. of Math. (2) \textbf{203} (2026), 623--675;
\href{https://arxiv.org/abs/2405.20552}{arXiv:2405.20552}.

\bibitem{HB}
D. R. Heath-Brown,
\emph{The differences between consecutive smooth numbers},
Acta Arith. \textbf{184} (2018), 267--285;
\href{https://arxiv.org/abs/1808.02947}{arXiv:1808.02947}.

\bibitem{HT}
A. Hildebrand and G. Tenenbaum,
\emph{Integers without large prime factors},
J. Th\'eor. Nombres Bordeaux \textbf{5} (1993), 411--484.

\bibitem{IK}
H. Iwaniec and E. Kowalski,
\emph{Analytic number theory},
Amer. Math. Soc. Colloq. Publ. \textbf{53}, 2004.

\bibitem{MR}
K. Matom\"aki and M. Radziwi\l\l,
\emph{Multiplicative functions in short intervals},
Ann. of Math. (2) \textbf{183} (2016), 1015--1056.

\bibitem{MT}
K. Matom\"aki and J. Ter\"av\"ainen,
\emph{Almost primes in almost all short intervals II},
Trans. Amer. Math. Soc. \textbf{376} (2023), 5433--5459;
\href{https://arxiv.org/abs/2207.05038}{arXiv:2207.05038}.

\bibitem{Ter}
J. Ter\"av\"ainen,
\emph{Almost primes in almost all short intervals},
Math. Proc. Cambridge Philos. Soc. \textbf{161} (2016), 247--281.

\end{thebibliography}

\end{document}
